%%% FIELD proof-rank-only-impossibility
Fix \(p\geq 2\) and a member of the one-root star family in \(\texttt{def:star-counterfamily}\).  Thus vertex \(1\) is the unique root and, for each \(j\in\{2,\ldots,p\}\), the only edge into \(j\) is \(1\to j\), with coefficient \(b_j\neq0\).  From \(A=I_p-B\) and \(a_j=A^{\top}e_j\), the transposes of the structural-equation rows are
\[
 a_1=e_1,
 \qquad
 a_j=e_j-b_je_1\quad (j=2,\ldots,p).
 \tag{1}
\]
All \(\omega_j\) and \(\tau_j\) are strictly positive.  Define
\[
 d=\tau_1^{-1}-\omega_1^{-1}.
\]
The star-family restriction gives \(d\neq0\); this is also the effectiveness condition for the root intervention because \(\operatorname{Pa}(1)=\varnothing\).  Applying \(\texttt{lem:surgery-identity}\) to the root and using (1) gives
\[
 \Delta_1=d e_1e_1^{\top}.
 \tag{2}
\]
For a leaf \(j\geq2\), the same lemma and (1) give
\[
 \Delta_j
 =\tau_j^{-1}e_je_j^{\top}
  -\omega_j^{-1}(e_j-b_je_1)(e_j-b_je_1)^{\top}.
 \tag{3}
\]
Therefore \(\Delta_j\) vanishes outside rows and columns \(\{1,j\}\), and its principal submatrix on those coordinates is
\[
 \left.\Delta_j\right|_{\{1,j\}}
 =
 \begin{pmatrix}
 -b_j^2/\omega_j & b_j/\omega_j\\
 b_j/\omega_j & \tau_j^{-1}-\omega_j^{-1}
 \end{pmatrix}.
 \tag{4}
\]
Direct expansion, using positivity of \(\omega_j\) and \(\tau_j\), yields
\[
 \begin{aligned}
 \det\!\left(\left.\Delta_j\right|_{\{1,j\}}\right)
 &= -\frac{b_j^2}{\omega_j}
       \left(\frac1{\tau_j}-\frac1{\omega_j}\right)
    -\frac{b_j^2}{\omega_j^2}\\
 &= -\frac{b_j^2}{\omega_j\tau_j}<0,
 \end{aligned}
 \tag{5}
\]
where strictness also uses \(b_j\neq0\).  Since all other rows and columns vanish, (5) proves
\[
 R_{0j}=\operatorname{rank}(K_j-K_0)
 =\operatorname{rank}(\Delta_j)=2.
 \tag{6}
\]
Subtracting (2) from (4) and expanding again gives
\[
 \begin{aligned}
 \det\!\left(
   \left.(\Delta_j-\Delta_1)\right|_{\{1,j\}}
 \right)
 &=\left(-\frac{b_j^2}{\omega_j}-d\right)
      \left(\frac1{\tau_j}-\frac1{\omega_j}\right)
      -\frac{b_j^2}{\omega_j^2}\\
 &=\frac{d(\tau_j-\omega_j)-b_j^2}{\omega_j\tau_j}.
 \end{aligned}
 \tag{7}
\]
The defining strict inequality in \(\texttt{def:star-counterfamily}\) makes (7) nonzero.  Hence
\[
 R_{1j}=\operatorname{rank}(K_j-K_1)
 =\operatorname{rank}(\Delta_j-\Delta_1)=2.
 \tag{8}
\]
Finally, (2) and \(d\neq0\) imply
\[
 R_{01}=\operatorname{rank}(K_1-K_0)=1.
 \tag{9}
\]

Let \(\sigma\) denote only the transposition of component labels \(0\) and \(1\), fixing every leaf-component label, and let \(P_\sigma\) be its permutation matrix on \(C\).  Equations (6) and (8) give \(R_{0j}=R_{1j}\) for every \(j\geq2\); equation (9) is symmetric in \(0\) and \(1\); entries indexed by two leaves are unchanged because \(\sigma\) fixes both indices; and every diagonal entry is zero.  These cases exhaust \(C\times C\), so
\[
 P_\sigma R P_\sigma^{\top}=R.
 \tag{10}
\]

The strict region is nonempty.  Set
\[
 \omega_k=1\quad(k\in V),\qquad
 \tau_1=2,\qquad
 \tau_j=1\quad(j=2,\ldots,p),\qquad
 b_j=1\quad(j=2,\ldots,p),
 \tag{11}
\]
and choose, for example, \(\pi_c=1/(p+1)\) for every \(c\in C\).  Then all variances and weights are positive, the weights sum to one, \(d=-1/2\neq0\), and for each leaf
\[
 d(\tau_j-\omega_j)-b_j^2=-1<0.
 \tag{12}
\]
The root intervention is effective because \(\tau_1\neq\omega_1\); every leaf intervention is effective because it has the parent \(1\), even though (11) has \(\tau_j=\omega_j\).  Together with the acyclic star graph and the baseline and surgery equations, these facts place the example in \(\mathcal M_p\) by \(\texttt{def:model-class}\).  Positivity, \(d\neq0\), the conditions \(b_j\neq0\), and the finitely many strict inequalities in (12) persist on a sufficiently small Euclidean neighborhood of the displayed variance and coefficient parameters; positivity of the weights persists on a relative neighborhood in the probability simplex.  Hence the family contains a nonempty relatively open strict-inequality parameter region.

For completeness, all \(p+1\) Gaussian components are distinct on this region.  Equation (2) separates the root intervention from the baseline; (5) separates every leaf intervention from the baseline; (7) separates every leaf intervention from the root intervention.  If \(j\neq k\) are leaves, (3) gives \((\Delta_j)_{1j}=b_j/\omega_j\neq0\), whereas \((\Delta_k)_{1j}=0\), so the two leaf components are distinct as well.

Now let \(F\) be any component-permutation-equivariant rank-only baseline selector.  By \(\texttt{def:rank-compression}\), it is a map from rank matrices to component labels satisfying
\[
 F(P_\gamma R P_\gamma^{\top})=\gamma(F(R))
 \tag{13}
\]
for every permutation \(\gamma\) of \(C\).  Apply (13) with \(\gamma=\sigma\) and use (10):
\[
 F(R)=\sigma(F(R)).
 \tag{14}
\]
Neither \(0\) nor \(1\) is fixed by \(\sigma\).  Thus (14) forces \(F(R)\in C\setminus\{0,1\}\); this remains valid for \(p=2\), when the fixed-point set is the singleton \(\{2\}\).  In particular, \(F(R)\neq0\), while component \(0\) is the observational baseline.  Therefore \(F\) fails at every point of the nonempty open region.  Since \(F\) was arbitrary, no component-permutation-equivariant rank-only selector universally recovers the baseline on \(\mathcal M_p\).

The same obstruction also transfers, negatively, to the full model class of Squires et al., without any positive recovery claim under latent mixing.  Indeed, \(\texttt{lem:squires-class-embedding}\) constructs a distinct reverse-topological latent-node bijection \(\rho\) and permutation matrix \(Qe_j=e_{\rho(j)}\), verifies the published assumptions, leaves the context labels unchanged, and proves \(\Theta_c=K_c\) for every \(c\in C\).  It therefore preserves every entry of \(R\).  The symbol \(\sigma\) in (10)--(14) remains solely the component-label transposition and is not the latent relabeling \(\rho\).  Consequently, any equivariant rank-only selector universally correct on the full Proposition-10 class would restrict to a selector correct on this embedded open star family, contradicting (14).
